\documentclass{article}

\usepackage{algorithmic}
\usepackage{xcolor}
\usepackage{sectsty}
\usepackage[algo2e, vlined,
			linesnumbered, ruled,
			algosection, nofillcomment]{algorithm2e}
\usepackage{setspace}

\title{\Huge \textbf{Typesetting Algorithms}}
\date{}

\sectionfont{\color{blue}}

\algsetup{indent = 2em,
		  linenosize = \large,
		  linenodelimiter = .
}

\renewcommand{\algorithmicrequire}{
	\color{red}{\textbf{Require:}}
}

\renewcommand{\algorithmicensure}{
	\color{blue}{\textbf{Ensure:}}
}

\SetKwInput{KwRequired}{Required}
\SetKwInput{KwExpected}{Expected}
\SetNlSty{texttt}{[}{]:}

\setstretch{1.1}

\begin{document}
	\maketitle
	
	\section{The algorithmic Package}
	\subsection{Conditional Statements and Looping}
	
	\begin{algorithmic}[1]
		\IF{Condition}
			\STATE Carry out some operations.
		\ELSIF{Another Condition}		
			\STATE Carry out another operation.
		\ELSE
			\STATE Do basic operations.
		\ENDIF
		% For Loop
		\FOR{$i=0$ \TO $10$}
			\STATE Carry out some operations.
		\ENDFOR
		% While Loop
		\WHILE{Some condition holds}
			\STATE Carry out some operations.
		\ENDWHILE
	\end{algorithmic}
	
	\bigskip
	\subsection{Conditions and Logical Operators}
	
	\begin{algorithmic}[1]
		% Pre-condition
		\REQUIRE $x \neq 0$ and $n \geq 0$
		% Post-condition
		\ENSURE $x \neq 0$ and $n \geq 0$
		% Logical Operators
		\IF{(value1 \AND value2) are equal}
			\STATE Do something 
		\ELSIF{\NOT(value1 $=$ value2)}
			\STATE Do something else.
		\ELSE[you can put a comment here]
			\RETURN \TRUE
		\ENDIF
	\end{algorithmic}
	
	\pagebreak
	\section{The algorithm2e Package}
	% Defining and Styling Functions
	\SetFuncSty{textbf}
	\SetKwFunction{intRest}{\color{blue}{IntervalRest}}
	\SetKwFunction{test}{\emph{TestFunc}}
	% Caption Styling
	\SetAlCapSty{textit}
	\SetAlCapNameSty{texttt}
	\SetAlCapFnt{\large}
	\SetAlCapNameFnt{\Large} 
	% Margins and Spaces
	\IncMargin{10pt}
	\SetAlgoInsideSkip{medskip}
	\setlength{\interspacetitleruled}{4pt}
	% Defining New Comments
	\SetKwComment{MatCmnt}{\% }{} 
	\begin{algorithm2e}[H]
		\DontPrintSemicolon
		%% Algorithm Head
		\KwRequired{This is a required stuff.}
		\KwExpected{This is the expected output range.}
		\KwData{$G=(X,U)$ such that $G^{tc}$ is an 					order.}
		\KwResult{$G'=(X,V)$ with $V \subseteq U$ such 				that $G'^{tc}$ is an interval order.}
		%% Algorithm Body
		\intRest{parameters} \\
		\test{parameters} 
		\MatCmnt*[h]{This is a Matlab Style Comment}\\
		$ V \longleftarrow U $ \;
		$ S \longleftarrow \emptyset $ \\
		% The First For Loop
		\For{$x \in X$}{
			$ NbSuccIns(x) \longleftarrow 0 $ \\
			$ NbPredInMin(x) \longleftarrow 0 $ \\
			$ NbPredNotInMin(x) \longleftarrow
				|ImPred(x)| $ 		
		}
		% The Second For Loop
		\For(\tcc*[h]{This is a C comment})
			{$x \in X$}{
			\If{
				$NbPredInMin(x)=0$ {\bf and}
				$NbPredNotInMin(x)=0$
			   }
			   {$AppendToMin(x)$}	
		}
		% The While Section
		\While{$S \neq \emptyset$}{
			% First Part
			remove $x$ from the list of $T$ of maximal 					index. \\		
			% The Second Part: The Nested While Loop
			\While{$ |S \cap ImSucc(x)| \neq |S| $}{
				 \For{$ y \in S-ImSucc(x) $}{
					\{ remove from $V$ all the				 					arcs $zy$: \} \\
					% The Nested For Loop
					\For{$z \in ImPred(y) \cap Min$}{
						% NFL Line No.1
						remove the arc $zy$ from $V$ \\
						% NFL Line No.2
						$NbSuccIns(z) \longleftarrow
							NbSuccIns(z)-1$ \\
						% NFL Line No.3
						move $z$ in $T$ to the list 							preceding its present list \\
						% NFL Line No.4
						\{i.e. If $z \in T[k]$,
							move $z$ from $T[k]$ to
							$T[k-1]$\}
						}
					% The Parent Loop Four Lines
					$NbPredInMin(y) \longleftarrow 0$\\
					$NbPredNotInMin(y) \longleftarrow 							0$ \\
					$S \longleftarrow S - \{y\}$\\
					$AppendToMin(y)$ \\
				 	}		  
				}
			% The Third Part
			$RemoveFromMin(x)$
		}		
		
		\caption{Interval Restriction}
	\end{algorithm2e}
\end{document}